% GENERATED FILE. Edit canonical lessons, then run npm run generate:books.
% canonical-source-hash: fnv1a64:9b76baea3c62f099
% canonical-lessons: BN-C79-bochhor, BN-C79-minit, BN-C79-aj, BN-C79-ekhon, BN-C79-sobsomoy

\chapter{Year, Minute, Today, Now, Always}
\label{ch:bn-year-minute-today-now-always}
\hlchaptermodality{79}

\begin{chapteropening}
\textbf{By the end of this chapter:} \emph{I can say a year, a minute, today, now and always.}

Complete the last lesson of chapter 79: I can say a year, a minute, today, now and always.
\end{chapteropening}

\section[\texorpdfstring{\bn{বছর} (bôchhôr)}{বছর (bôchhôr)}]{\bn{বছর} (bôchhôr) — a year}

\label{lesson:BN-C79-bochhor}



\begin{quote}
Before the new one: say the Bengali for a week, then the Bengali for a month.
\end{quote}

\subsection*{You'll want to know: \bn{বছর}}
\textbf{\bn{বছর}} — \emph{bôchhôr} — \textquotedblleft{}a year\textquotedblright{}.

\begin{cousinweb}
\textbf{\bn{বছর}} (\emph{bôchhôr}) goes back to Sanskrit \emph{vatsara}, \textquotedblleft{}year\textquotedblright{} (CDIAL 11242). CDIAL asks whether the Bengali form was borrowed from Sanskrit rather than inherited.
\end{cousinweb}

\begin{grammarlens}[title={What you've built}]
One more word for this chapter.
\end{grammarlens}

\subsection*{Your turn}
\begin{itemize}
\raggedright
  \item \emph{Say it:} \emph{bôchhôr}
  \item \emph{Say it:} \emph{bôchhôr}, once more
  \item \emph{Say it:} say \emph{bôchhôr}
  \item \emph{Recall:} say the Bengali for a week, then the Bengali for a month, then say \emph{bôchhôr} again
\end{itemize}

\begin{tcolorbox}[breakable,colback=teal!4,colframe=teal!35!black,title={Before you move on}]
What does \bn{বছর} mean? (\textquotedblleft{}A year\textquotedblright{}.) Say it once more.
\end{tcolorbox}
\section[\texorpdfstring{\bn{মিনিট} (miniṭ)}{মিনিট (miniṭ)}]{\bn{মিনিট} (miniṭ) — a minute}

\label{lesson:BN-C79-minit}



\begin{quote}
Before the new one: say the Bengali for a month, then the Bengali for a year.
\end{quote}

\subsection*{You'll want to know: \bn{মিনিট}}
\textbf{\bn{মিনিট}} — \emph{miniṭ} — \textquotedblleft{}a minute\textquotedblright{}.

\begin{cousinweb}
\textbf{\bn{মিনিট}} (\emph{miniṭ}) is a loan from English \emph{minute}.
\end{cousinweb}

\begin{grammarlens}[title={What you've built}]
One more word for this chapter.
\end{grammarlens}

\subsection*{Your turn}
\begin{itemize}
\raggedright
  \item \emph{Say it:} \emph{miniṭ}
  \item \emph{Say it:} \emph{miniṭ}, once more
  \item \emph{Say it:} say \emph{miniṭ}
  \item \emph{Recall:} say the Bengali for a month, then the Bengali for a year, then say \emph{miniṭ} again
\end{itemize}

\begin{tcolorbox}[breakable,colback=teal!4,colframe=teal!35!black,title={Before you move on}]
What does \bn{মিনিট} mean? (\textquotedblleft{}A minute\textquotedblright{}.) Say it once more.
\end{tcolorbox}
\section[\texorpdfstring{\bn{আজ} (āj)}{আজ (āj)}]{\bn{আজ} (āj) — today}

\label{lesson:BN-C79-aj}



\begin{quote}
Before the new one: say the Bengali for a year, then the Bengali for a minute.
\end{quote}

\subsection*{You'll want to know: \bn{আজ}}
\textbf{\bn{আজ}} — \emph{āj} — \textquotedblleft{}today\textquotedblright{}.

\begin{cousinweb}
\textbf{\bn{আজ}} (\emph{āj}) is a tadbhava, a word inherited from Sanskrit and reshaped by sound change. It comes from \emph{adya}, \textquotedblleft{}today\textquotedblright{} (CDIAL 242).
\end{cousinweb}

\begin{grammarlens}[title={What you've built}]
One more word for this chapter.
\end{grammarlens}

\subsection*{Your turn}
\begin{itemize}
\raggedright
  \item \emph{Say it:} \emph{āj}
  \item \emph{Say it:} \emph{āj}, once more
  \item \emph{Say it:} say \emph{āj}
  \item \emph{Recall:} say the Bengali for a year, then the Bengali for a minute, then say \emph{āj} again
\end{itemize}

\begin{tcolorbox}[breakable,colback=teal!4,colframe=teal!35!black,title={Before you move on}]
What does \bn{আজ} mean? (\textquotedblleft{}Today\textquotedblright{}.) Say it once more.
\end{tcolorbox}
\section[\texorpdfstring{\bn{এখন} (ekhôn)}{এখন (ekhôn)}]{\bn{এখন} (ekhôn) — now}

\label{lesson:BN-C79-ekhon}



\begin{quote}
Before the new one: say the Bengali for a minute, then the Bengali for today.
\end{quote}

\subsection*{You'll want to know: \bn{এখন}}
\textbf{\bn{এখন}} — \emph{ekhôn} — \textquotedblleft{}now\textquotedblright{}.

\begin{cousinweb}
\textbf{\bn{এখন}} (\emph{ekhôn}) is the near head \emph{e-} on \emph{-khôn}, from Sanskrit \emph{kṣaṇa}, \textquotedblleft{}moment\textquotedblright{}: CDIAL lists \emph{takhan}, \textquotedblleft{}then\textquotedblright{}, under Sanskrit \emph{tatkṣaṇa}, \textquotedblleft{}that very moment\textquotedblright{}, with the same ending on the other pronoun stems (CDIAL 5640).
\end{cousinweb}

\begin{grammarlens}[title={What you've built}]
One more word for this chapter.
\end{grammarlens}

\subsection*{Your turn}
\begin{itemize}
\raggedright
  \item \emph{Say it:} \emph{ekhôn}
  \item \emph{Say it:} \emph{ekhôn}, once more
  \item \emph{Say it:} say \emph{ekhôn}
  \item \emph{Recall:} say the Bengali for a minute, then the Bengali for today, then say \emph{ekhôn} again
\end{itemize}

\begin{tcolorbox}[breakable,colback=teal!4,colframe=teal!35!black,title={Before you move on}]
What does \bn{এখন} mean? (\textquotedblleft{}Now\textquotedblright{}.) Say it once more.
\end{tcolorbox}
\section[\texorpdfstring{\bn{সবসময়} (sôbsômôy)}{সবসময় (sôbsômôy)}]{\bn{সবসময়} (sôbsômôy) — always}

\label{lesson:BN-C79-sobsomoy}



\begin{quote}
Before the new one: say the Bengali for today, then the Bengali for now.
\end{quote}

\subsection*{You'll want to know: \bn{সবসময়}}
\textbf{\bn{সবসময়}} — \emph{sôbsômôy} — \textquotedblleft{}always\textquotedblright{}.

\begin{cousinweb}
\textbf{\bn{সবসময়}} (\emph{sôbsômôy}) is \emph{sab}, \textquotedblleft{}all\textquotedblright{}, from Sanskrit \emph{sarva}, \textquotedblleft{}all\textquotedblright{} (CDIAL 13276), + \emph{sômôy}, \textquotedblleft{}time\textquotedblright{}, Sanskrit \emph{samaya} taken over whole.
\end{cousinweb}

\begin{grammarlens}[title={What you've built}]
One more word for this chapter.
\end{grammarlens}

\subsection*{Your turn}
\begin{itemize}
\raggedright
  \item \emph{Say it:} \emph{sôbsômôy}
  \item \emph{Say it:} \emph{sôbsômôy}, once more
  \item \emph{Say it:} say \emph{sôbsômôy}
  \item \emph{Recall:} say the Bengali for today, then the Bengali for now, then say \emph{sôbsômôy} again
\end{itemize}

\begin{tcolorbox}[breakable,colback=teal!4,colframe=teal!35!black,title={Before you move on}]
What does \bn{সবসময়} mean? (\textquotedblleft{}Always\textquotedblright{}.) Say it once more.
\end{tcolorbox}
